\documentclass[10pt,a4paper]{amsart}
\usepackage[margin=19mm]{geometry}
\usepackage{amsmath,amssymb,mathtools}
\usepackage[scr=rsfs,cal=euler]{mathalfa}
\usepackage{xfrac}
\usepackage[unicode,hidelinks,bookmarks=false]{hyperref}
\newcommand{\ie}{i.e.\ }
\newcommand{\cf}{cf.\ }
\newcommand{\resp}{respectively\ }
\newcommand{\Subsection}{\subsection}
\providecommand{\nobreakdash}{\nobreak\hbox{-}}
\setlength{\emergencystretch}{2em}
\multlinegap0pt
\usepackage[shortlabels]{enumitem}
\usepackage{amssymb}
\usepackage{tikz}
\usepackage{mathtools}
\newtheorem{lemma}{Lemma}
\newcommand{\R}{\mathbb R}
\title{Volume and Projection Inequalities II:\\ Determinants and $L_p$-Sums}
\author{Matthieu Fradelizi, Auttawich Manui, Cheikh Saliou Ndiaye, Artem Zvavitch}
\date{Source: arXiv:2608.24081v1 (CC BY 4.0)}
\begin{document}
\maketitle

\noindent\emph{Source and licence.} Volume and Projection Inequalities II:\\ Determinants and $L_p$-Sums. Version arXiv:2608.24081v1 (CC BY 4.0), distributed by arXiv under the Creative Commons Attribution 4.0 International licence, http://creativecommons.org/licenses/by/4.0/, as recorded in the arXiv licence metadata for this exact version and shown on the abstract page https://arxiv.org/abs/2608.24081v1. Full licence notice: \texttt{licenses/arxiv-2608.24081-CC-BY-4.0.txt}.\par\smallskip

\noindent\textit{Editorial scope.} Counted unit: the article's lemma lem:mazur-map-section (source0.tex lines 734--750). The article's complete original proof of it is delivered with it (source0.tex lines 1846--2075, one \texttt{proof} environment of 230 lines, which the article places away from the statement). No mathematics is written, repaired or re-derived: the statement and the proof are the article's own lines, in the article's own notation, and no retained line is substituted or re-flowed.  No further source-local context is needed: every cross-reference of every retained line resolves inside the two delivered spans.  Nothing was omitted from any retained span, so the package carries no omission record.  No retained line contains a citation, so the wrapper carries no bibliography and no bibliography entry.  No retained line contains a figure environment, an image-inclusion command or drawing code, so no image is delivered and the package declares no asset dependency.  Editorial formatting only, in addition: an AMS \texttt{amsart} preamble that restates the article's own declarations and macros used by the retained lines, namely the environments \texttt{lemma} on separate counters, together with the standard packages those lines need.  The retained environments print in this document as Lemma 1 in order of appearance, whereas the article prints the same items with the numbering of its own full document.  The complete native source of the article as distributed in the arXiv e-print for this exact version is bundled unchanged as \texttt{sources/208447/source0.tex}.
\begin{lemma}
\label{lem:mazur-map-section}
Let $1<p\leq2$ and let $q$ be  the H\"older conjugate of $p$, that is $\frac{1}{p} + \frac{1}{q} = 1$. For a nonzero vector $u\in\R^m$, define
\[
    J(u)=\frac{(|u_1|^{p-2}u_1,\ldots,|u_m|^{p-2}u_m)}{\|u\|_p^{p-1}}
\]
to be the vector in $\partial B_q^m$ such that $\langle J(u),u\rangle=\|u\|_p$.
Then, for nonzero $u,v\in\R^m$,
\begin{equation}
\label{eq:mazur-map-estimate-section}
    \|J(u)-J(v)\|_q
    \leq
    2\left(
       \frac{\|u-v\|_p}{\|u\|_p+\|v\|_p}
    \right)^{p-1}.
\end{equation}
\end{lemma}

\begin{proof}[Proof of Lemma~\ref{lem:mazur-map-section}] 
Set $\phi(t)=|t|^{p-2}t,$ for $ t\in\R\setminus\{0\},$ and $\phi(0)=0.$ We first prove the following scalar inequality. If
$\lambda,\mu\geq0$, $\lambda+\mu=1$, and $s,t\in\R$, then
\begin{equation}
\label{eq:scalar-mazur-map-section-claim}
    |\phi(s)-\phi(t)|^q
    \leq
    2^q|\lambda s-\mu t|^p
    +
    2^{q-1}(\mu-\lambda)
    \bigl(|s|^p-|t|^p\bigr).
\end{equation}
The inequality is invariant under the simultaneous interchange $(s,\lambda)\leftrightarrow(t,\mu)$. Thus, we may assume that $|s| \geq |t|$. By continuity, it suffices to consider the case $
|s|\geq |t|>0$. The inequality is also unchanged if $(s,t)$ is replaced by $(-s,-t)$, so we may assume that $t > 0$. Since the inequality is homogeneous of degree $p$, we divide by $t$ and reduce to the case $t=1$. Hence, by substituting $\mu = 1-\lambda$, it remains to prove for $s \in \R$ with $|s| \geq 1$, and $\lambda \in [0,1]$,
\begin{equation}
\label{eq:scalar-mazur-map-section}
    |\phi(s)-1|^q
    \leq
    2^q|\lambda (s+1)-1 |^p
    +
    2^{q-1}(1-2\lambda)
    \bigl(|s|^p-1\bigr).
\end{equation}
\begin{enumerate}[wide =0pt, labelwidth= 1.1cm]
    \item[\bf Case 1:] Suppose that $s \geq 1$. If $s = 1$, then inequality \eqref{eq:scalar-mazur-map-section} is immediate since the left-hand side is zero. Hence, we assume that $s > 1$. Fix $s$. We minimize the right-hand side of \eqref{eq:scalar-mazur-map-section} with respect to $\lambda$. Define a function $F$ on $[0,1]$ by 
    \[
        F(\lambda) = 2^q|\lambda(s+1)-1|^p+2^{q-1}(1-2 \lambda)\left(s^p-1\right).
    \]
    Differentiating $F$ we obtain 
    \[
        F^{\prime}(\lambda)=2^q p(s+1) \operatorname{sgn}(\lambda(s+1)-1)|\lambda(s+1)-1|^{p-1}-2^q\left(s^p-1\right).
    \]
    The critical-point equation is 
    $$
        p(s+1)\operatorname{sgn}(\lambda(s+1)-1)|\lambda(s+1)-1|^{p-1}=s^p-1 .
    $$
    Since $ s > 1$, the right-hand side is positive. Therefore,  any solution must satisfy $\lambda > \frac{1}{s+1}$. Solving the equation, the critical point is
    \[
        \lambda_0 =\frac{1+\left(\frac{s^p-1}{p(s+1)}\right)^{1 /(p-1)}}{s+1}.
    \]
    Since $1< p \leq 2$, the function $x \mapsto x^{p-1}$ is concave. Using Jensen's inequality, we have 
    \begin{equation}
    \label{eq:Jen-with-s}
        \frac{s^p-1}{p(s-1)} =\frac{1}{s-1}\int_1^s t^{p-1} dt \leq \left(\frac{1}{s-1} \int_1^s t d t \right)^{p-1} = \left(\frac{s+1}{2}\right)^{p-1}.
    \end{equation}
    Hence,
    \begin{equation} 
    \frac{s^p-1}{p(s+1)}
\leq
    \frac{s-1}{s+1}
    \left(\frac{s+1}{2}\right)^{p-1}
    \leq
    \left(\frac{s-1}{2}\right)^{p-1}. \label{eq:bound-of-F}
    \end{equation}
    Therefore, $\lambda_0 \leq \frac{1}{2}$. Since $F$ is strictly convex, $\lambda_0$ is the only minimizer of $F$. 
    Substituting $\lambda_0$ into $F$, we obtain 
    \begin{align}
        F\left(\lambda_0\right)&=2^q  \left(\frac{s^p-1}{p(s+1)}\right)^{p /(p-1)} 
         +2^{q-1}\left(1-\frac{2+2\left(\frac{s^p-1}{p(s+1)}\right)^{1 /(p-1)}}{s+1}\right)\left(s^p-1\right)
         % \\
         % & =\frac{2^q\left(s^p-1\right)}{p(s+1)}\left(\frac{s^p-1}{p(s+1)}\right)^{1 /(p-1)} +\frac{2^{q-1}\left(s^p-1\right)}{s+1}\left[s-1-2\left(\frac{s^p-1}{p(s+1)}\right)^{1 /(p-1)}\right]
         % \\
         % &=\frac{2^{q-1}\left(s^p-1\right)}{s+1}\left[\frac{2}{p}\left(\frac{s^p-1}{p(s+1)}\right)^{1 /(p-1)}+s-1-2\left(\frac{s^p-1}{p(s+1)}\right)^{1 /(p-1)}\right]
         \\ & 
         =\frac{2^{q-1}\left(s^p-1\right)}{s+1}\left[s-1-\frac{2(p-1)}{p}\left(\frac{s^p-1}{p(s+1)}\right)^{1 /(p-1)}\right] .
    \end{align}
    Using \eqref{eq:bound-of-F}, we obtain
    \[
        F(\lambda_0) \geq \frac{2^{q-1}\left(s^p-1\right)}{s+1} \cdot \frac{s-1}{p}.
    \]
    To verify \eqref{eq:scalar-mazur-map-section}, it is enough to prove
    \begin{equation} \label{eq:to-prove-final-s}
        |s^{p-1} -1|^q \leq \frac{2^{q-1}\left(s^p-1\right)(s-1)}{p(s+1)}, \qquad s> 1.
    \end{equation}
    Define $G$ on $[1,\infty)$ by
    \[
        G(s)=
        \frac{
            2^{q-1}(s^p-1)(s-1)
        }{
            p(s+1)(s^{p-1}-1)^q
        }.
    \]
    Taking logarithms on both sides and then differentiating, we obtain
    \[
        \frac{G'(s)}{G(s)}
   =
    \frac{p s^{p-1}}{s^p-1}
    +
    \frac{1}{s-1}
    -
    \frac{1}{s+1}
    -
    \frac{p s^{p-2}}{s^{p-1}-1}
=
    -\frac{
        p s^{p-2}(s-1)
    }{
        (s^p-1)(s^{p-1}-1)
    }
    +
    \frac{2}{s^2-1}.
    \]
    Claim that $G' (s) \leq 0$. Equivalently, 
    \[
        2 \left( \frac{(s^p -1)(s^{p-1} -1)}{p(s-1)^2} \right) 
        \leq s^{p-2} (s+1).
    \]
    Using \eqref{eq:Jen-with-s} and the Hermite--Hadamard inequality with the convex function $t \mapsto t^{p-2}$ on $ [1,s]$, we obtain
    \begin{align}
        2 \left( \frac{(s^p -1)(s^{p-1} -1)}{p(s-1)^2} \right) &=2 (p-1)\left( \frac{s^p -1}{p(s-1)} \right) \left( \frac{s^{p-1} -1}{(p-1)(s-1)} \right)
        \\
        &\leq 
        2(p-1) \left(\frac{s+1}{2}\right)^{p-1}\left( \frac{1}{s-1} \int_{1}^s t^{p-2} dt \right)
        \\
        &\leq 
        2(p-1) \left(\frac{s+1}{2}\right)^{p-1} \left( \frac{s^{p-2}+1}{2} \right)
        \\
        &\leq 2^{1-p}s^{p-2}(s+1)^{p-1} (1 + s^{2-p}). \label{eq:final-first-case}
    \end{align}
    Using H\"older's inequality, we obtain
    \[
        1 + s^{2-p} \leq  2^{p-1}(1+s)^{2-p}.
    \]
    Plugging it into \eqref{eq:final-first-case}, we verify that $G'(s) \leq 0$ for $s >1$. Moreover, since $ q \geq 2$ and $p \leq 2$, we get
\[
    \lim_{s\to\infty}G(s)
    =
    \frac{2^{q-1}}{p} \geq 1.
\]
Thus, $G(s)\geq1$ for $s >1$, proving \eqref{eq:to-prove-final-s} and hence
\eqref{eq:scalar-mazur-map-section} in the case $s\geq1$.
\item[\bf Case 2:] Suppose that $s \leq -1$. If $s = -1$, then both sides of inequality \eqref{eq:scalar-mazur-map-section} are equal. Thus, we only need to prove when $s < -1$. To avoid confusion caused by the sign of $s$, we write $s=-r$ where $ r > 1$. Thus, we need to prove
\begin{equation}
\label{eq:scalar-mazur-map-section-2}
    (r^{p-1}+1)^q
    \leq
    2^q(1+ \lambda (r-1))^p
    +
    2^{q-1}(1-2\lambda)
    \bigl(r^p-1\bigr).
\end{equation}
Fix $r$, and define $F$ on $[0,1]$ by
$$
F(\lambda)=2^q(1+\lambda(r-1))^p+2^{q-1}(1-2 \lambda)\left(r^p-1\right) .
$$
Differentiating, we obtain
$$
F^{\prime}(\lambda)=2^q p(r-1)(1+\lambda(r-1))^{p-1}-2^q\left(r^p-1\right) .
$$
The critical-point equation is
$
p(r-1)(1+\lambda(r-1))^{p-1}=r^p-1 .
$
Therefore, the critical point is
$$
\lambda_0=\frac{\left(\frac{r^p-1}{p(r-1)}\right)^{1 /(p-1)}-1}{r-1} .
$$
Since $t^{p-1} \geq 1$ on $[1, r]$,
$$
\frac{r^p-1}{p(r-1)}=\frac{1}{r-1} \int_1^r t^{p-1} d t \geq 1.
$$
Hence,
$
\lambda_0 \geq 0 .
$
Moreover, since $1<p \leq 2$, the function $t \mapsto t^{p-1}$ is concave. By Jensen's inequality,
\begin{equation} \label{eq:second-case-claim-001}
\frac{r^p-1}{p(r-1)} \leq \left(\frac{r+1}{2}\right)^{p-1}.
\end{equation}
Thus, $\lambda_0 \leq \frac{1}{2}$. Since $F$ is strictly convex, $\lambda_0$ is its unique minimizer. Substituting $\lambda_0$ into $F$, we obtain
\[
    F(\lambda_0 ) = 2^{q-1} \frac{r^p-1}{p(r-1)}\left[p(r+1)-2(p-1)\left(\frac{r^p-1}{p(r-1)}\right)^{1 /(p-1)}\right].
\]
Using \eqref{eq:second-case-claim-001}, we obtain
$$
F\left(\lambda_0\right) \geq 2^{q-1} \frac{r^p-1}{p(r-1)}(r+1) .
$$
We use the Hermite--Hadamard inequality. Since $t \mapsto t^{p-1}$ is concave,
$$
\frac{r^p-1}{p(r-1)}=\frac{1}{r-1} \int_1^r t^{p-1} d t \geq \frac{1+r^{p-1}}{2}.
$$
It follows that
$
F\left(\lambda_0\right) \geq 2^{q-2}\left(1+r^{p-1}\right)(r+1) .
$
To prove \eqref{eq:scalar-mazur-map-section-2}, it remains to check
\begin{equation} \label{eq:final-case-2}
\left(r^{p-1}+1\right)^q \leq 2^{q-2}\left(1+r^{p-1}\right)(r+1).
\end{equation}
By H\"older's inequality, we have
$
1+r^{p-1} \leq 2^{2-p}(1+r)^{p-1} .
$
Thus,
$$
    \left(1+r^{p-1}\right)^{q-1}  \leq 2^{(2-p)(q-1)}(r+1)^{(p-1)(q-1)}  =2^{q-2}(r+1) ,
    $$
    proving \eqref{eq:final-case-2}, hence \eqref{eq:scalar-mazur-map-section-2}.
\end{enumerate}

    Combining these two cases proves \eqref{eq:scalar-mazur-map-section-claim}. 
    Write
    \[
    J(u)=\left(\phi\left(\frac{u_1}{\|u\|_p}\right), \ldots, \phi\left(\frac{u_m}{\|u\|_p}\right)\right),
    \quad \text{and}\quad
    J(v)=\left(\phi\left(\frac{v_1}{\|v\|_p}\right), \ldots, \phi\left(\frac{v_m}{\|v\|_p}\right)\right).
\]
Applying \eqref{eq:scalar-mazur-map-section-claim} with
$s=\frac{u_i}{\|u\|_p}$, $t=\frac{v_i}{\|v\|_p}$ and $\lambda = \frac{\|u\|_p}{\|u\|_p + \|v\|_p}$, we obtain
\[
\left|\phi\left(\frac{u_i}{\|u\|_p}\right)-\phi\left(\frac{v_i}{\|v\|_p}\right)\right|^q
    \leq
    2^q\frac{|u_i -v_i|^p}{(\|u\|_p + \|v\|_p)^p}
    +
    2^{q-1}\left(\frac{\|v\|_p - \|u\|_p}{\|u\|_p + \|v\|_p}\right)
    \left(\frac{\left|u_i\right|^p}{\|u\|_p^p}-\frac{\left|v_i\right|^p}{\|v\|_p^p}\right).
\]
Summing over $i$, we have
\[
    \|J(u)-J(v)\|_q^q
    \leq
    2^q
    \left(
        \frac{\|u-v\|_p}
        {\|u\|_p+\|v\|_p}
    \right)^p,
\]
where the last term vanishes. Therefore, raising both sides to the power $1 / q$ completes the proof.
\end{proof}

\end{document}
